\section{Method}
\label{sec:method}

\subsection{Pointwise baseline}

For each nuclide and library temperature, $\sigma(E_i, T_j)$
lies on a union energy grid $\{E_i\}_{i=0}^{N_E - 1}$. For
$E_i \leq E < E_{i+1}$ the standard log-log lookup is
\begin{equation}
\sigma(E) = \sigma_i
\left(\sigma_{i+1}/\sigma_i\right)^{f},
\qquad
f = \frac{\ln(E/E_i)}{\ln(E_{i+1}/E_i)}.
\label{eq:loglog}
\end{equation}
We use a log-uniform hash index with $8192$ bins
(Brown 2014~\cite{brown2014hash}) to resolve $i$ in $O(1)$,
falling back to a bounded binary search in dense bins.
Non-positive cross-section values fall back to linear
interpolation to avoid $\log$ of a negative.

\subsection{Truncated SVD}

Pre-clamp the matrix to avoid $-\infty$ in the logarithm,
$\mathbf{A}_{ij} = \max(\sigma(E_i, T_j), 10^{-30})$, and
compute the thin SVD of $\mathbf{L} = \log_{10}\mathbf{A}$:
\begin{equation}
\mathbf{L} = \mathbf{U} \boldsymbol{\Sigma} \mathbf{V}^\top,
\qquad
\mathbf{L} \approx \mathbf{U}_k \boldsymbol{\Sigma}_k
\mathbf{V}_k^\top.
\end{equation}
Pre-multiply the basis $\mathbf{B} = \mathbf{U}_k
\boldsymbol{\Sigma}_k$ so a reconstruction at grid point $i$
for target coefficient vector $\mathbf{c} \in \mathbb{R}^k$ is
one dot product:
\begin{equation}
\log_{10}\tilde\sigma(E_i, T)
= \sum_{\ell = 1}^{k} B_{i\ell}\, c_\ell,
\qquad
\tilde\sigma = 2^{(\log_{10}\tilde\sigma)\cdot\log_2 10}.
\label{eq:svd_eval}
\end{equation}
$\mathbf{B}$ stores $N_E k$ doubles row-major so $k$ values
for one energy point live in $8k$ contiguous bytes; at $k=5$
that is one cache line. Between grid points we apply log-log
interpolation on $\log_{10}\tilde\sigma$ and take one
$\mathrm{exp}_2$.

\subsection{Temperature interpolation}
\label{sec:ducru}

Library temperatures are typically $\{0, 250, 294, 600, 900,
1200, 2500\}$~K in the ENDF/B-VII.1 HDF5 distribution. For
target $T \not\in \{T_j\}$ the coefficient vector follows
the Ducru kernel method~\cite{ducru2017}:
\begin{equation}
c_\ell(T) = \sum_{j \in \mathcal{S}} w_j(T)\, V_{j\ell},
\qquad
w_j(T) = \frac{\sqrt{T_j T}}{T_j + T}
\prod_{i \in \mathcal{S}, i \neq j}
\frac{T - T_i}{T + T_i} \cdot \frac{T_j + T_i}{T_j - T_i},
\label{eq:ducru}
\end{equation}
where $\mathcal{S}$ is a subset of library temperatures used
for the interpolation. Eq.~\ref{eq:ducru} is exact at library
temperatures ($w$ collapses to a one-hot selector when
$T = T_j$), $\sim\!0.1\%$ error on Doppler-broadened cross
sections over $300$--$3000$~K, and L2-optimal for the free
Doppler kernel on $\mathcal{S}$.

\subsubsection*{Choosing the interpolation subset}

Using $|\mathcal{S}| = N$ (all library temperatures) is
numerically unstable for $N \gtrsim 4$ because the
product-of-ratios term in $w_j$ develops near-singular
factors when non-adjacent $T_i$'s bracket $T$ with large
stride. In our measurements with $N = 7$ the $2500$~K column
pulls the product toward values $\sim\!10^{2}$, the resulting
weights routinely swing by a factor of $10$--$100$ across
neighbouring target temperatures, and the log-space
reconstruction sees non-physical oscillations that scale
linearly with the condition number.

We use the three library temperatures nearest to $T$.
Three is the smallest subset that captures a quadratic
correction to the Faddeeva kernel (two-point misses it; see
Sec.~\ref{sec:off_library}); it is also small enough that
the product term has only two factors per weight and the
formula is well-conditioned at every operating temperature
we tested in the interval $[250, 2500]$~K. Section~\ref{sec:ducru_puoi}
describes how we make the subset weights a partition of
unity; this is essential for resonance-dominated systems
and is the change that brings the off-library PWR result
into the deployable band against ACE+WMP.

\subsection{Partition-of-unity Ducru for off-library \texorpdfstring{$T$}{T}}
\label{sec:ducru_puoi}

When Eq.~\ref{eq:ducru} is applied to $V^\top$ columns of a
$\log_{10}\sigma$ SVD, the reconstructed
$\log_{10}\tilde\sigma(E_i, T)$ is a linear combination of
the per-library-$T$ log cross sections with coefficients
$w_j(T)$. On the \emph{linear} cross section this is a
weighted geometric mean:
\begin{equation}
\tilde\sigma(E, T)
= \prod_{j \in \mathcal{S}}
   \sigma(E, T_j)^{w_j(T)}.
\label{eq:geom_mean}
\end{equation}
If $\sum_j w_j \neq 1$, the right-hand side of
Eq.~\ref{eq:geom_mean} is scaled uniformly by
$\exp[(\sum_j w_j - 1) \cdot \overline{\log\sigma}]$ where
the overbar denotes a weighted average of
$\log_{10}\sigma(E, T_j)$. For smooth channels this is a
barely-visible gain error (far away from resonance peaks,
the average $\log\sigma$ is small and the multiplicative
shift is close to unity). On narrow resonance peaks it is
catastrophic: at the U-$238$ $6.67$~eV resonance with a
peak cross section of $\sim\!4500$~b, a $5\%$ gain error
moves the peak by $\sim\!225$~b, which on a PWR pin cell
translates to $\sim\!2000$~pcm on $k_\infty$
(Sec.~\ref{sec:off_library}).

Raw Ducru does \emph{not} have $\sum_j w_j = 1$ in general.
Eq.~\ref{eq:ducru}'s weights are designed to reconstruct
$\sigma(E, T)$ in the free-Doppler-kernel L2 norm
(see~\cite{ducru2017}, §3); partition of unity is not one
of the optimization's constraints. Measured on our
$N = 3$ bracketing subset at $T = 750$~K with library
$(600, 900, 1200)$~K the raw weights
$(0.21, 0.30, -0.01)$ sum to $0.50$. At $T = 1050$~K with
library $(900, 1200, 2500)$~K they sum to $0.51$. At
endpoints they collapse to $(1,0,0)$ / $(0,1,0)$ so the
total is $1$; \emph{between} endpoints the sum drifts by
up to $\approx\!0.5$.

We restore partition of unity by a single division:
\begin{equation}
\tilde w_j(T) = \frac{w_j(T)}{\sum_{k \in \mathcal{S}} w_k(T)},
\qquad
\sum_{j \in \mathcal{S}} \tilde w_j(T) = 1.
\label{eq:ducru_unity}
\end{equation}
This is the scheme used throughout the rest of the paper
for off-library $T$. It preserves the Faddeeva-derived
\emph{ratio} $w_j / w_k$ (the physically informative piece
of Eq.~\ref{eq:ducru}), absorbs the scalar magnitude error
into a pure rescaling, and leaves Eq.~\ref{eq:ducru_unity}
exact at library endpoints because the one-hot pattern
survives normalization. Sec.~\ref{sec:off_library} measures
the $k_\infty$ consequence on PWR and Godiva.

\subsubsection*{What $\tilde w_j$ is \emph{not}}

It is not the Ducru QP from~\cite{ducru2017} §4. The QP
enforces $\sum_j w_j = 1$ \emph{and} $w_j \geq 0$
\emph{and} optimizes the L2 fit in the Doppler-kernel
norm. Eq.~\ref{eq:ducru_unity} enforces only the first
constraint and does so as a rescaling rather than a
projection, so the result is still signed
($\tilde w_j$ can be negative). On U-$238$~MT$=102$ held
out at $900$~K, the non-negative simplex
projection~\cite{duchi2008simplex} of the raw Ducru
weights raises the L2 reconstruction error from $1.5\%$ to
$5\%$ relative to the signed scheme
(Sec.~\ref{sec:off_library_qp}). For the benchmarks
reported here, Eq.~\ref{eq:ducru_unity} suffices; an
analytic kernel-QP is left to future work.

\subsection{Memory trade-off of truncated SVD}
\label{sec:svd_tradeoff}

Truncated SVD is not a free memory win, and the comparison
against a pointwise table depends on how many temperatures
the library stores. A single-temperature pointwise table
stores one value per energy point per reaction. An SVD basis
at rank $k$ stores $k$ values per energy point per reaction
regardless of the number of library temperatures it was
trained on. Writing $T$ for the stored temperature count:

\begin{equation}
\frac{\text{SVD bytes}}{\text{pointwise bytes}}
= \frac{N_E \cdot k}{N_E \cdot T}
= \frac{k}{T}.
\label{eq:svd_pt_ratio}
\end{equation}

So:
\begin{itemize}
\item At $k = 1$, SVD is the same size as a
  single-temperature pointwise table but represents all $T$
  temperatures; SVD wins on memory against any pointwise
  table with $T \geq 1$ stored temperatures, and does
  strictly better as $T$ grows.
\item At $k \geq 2$, SVD only beats pointwise when the
  library carries $T > k$ temperatures. Break-even sits
  around $T = 2$--$3$. For a single-temperature toy
  problem SVD is strictly larger.
\item Production libraries typically carry $T \in [6, 15]$
  temperatures (depletion feedback, fuel-temperature
  reactivity), so the rank-$5$ SVD used in this paper
  ($k/T = 5/6 = 0.83$ to $5/15 = 0.33$) is a modest to
  strong memory win in the regime that matters, and a loss
  in the regime that does not.
\end{itemize}

The cost side of that trade-off is not zero:

\begin{itemize}
\item \textbf{Compute per lookup.} Pointwise is a hashed
  index plus one log-log interpolation. SVD is a rank-$k$
  FMA sequence plus a Ducru temperature step
  (Sec.~\ref{sec:ducru}). At $k = 5$ on the large-L3 system, SVD is
  $43$~ns/lookup against $91$~ns for the table
  (Sec.~\ref{sec:xs_pareto}) --- but this advantage depends
  entirely on the basis fitting in L2/L3 cache. A
  higher-rank basis or a larger nuclide set could tip this
  back toward pointwise.
\item \textbf{Reconstruction error.} Pointwise is exact at
  grid points and piecewise-linear-in-log between them; SVD
  is an approximation everywhere. RMSE on $\log_{10}\sigma$
  is $6\cdot 10^{-4}$ at $k = 5$ and reaches machine
  precision at $k = 6$ (Sec.~\ref{sec:xs_pareto}). Lower
  rank means smaller memory and faster lookup at the cost
  of a tunable amount of cross-section error, which
  propagates into $k_\text{eff}$ as pcm-scale bias.
\item \textbf{One-time decomposition cost.} SVD requires an
  offline factorisation per nuclide per reaction (seconds at
  load time). Pointwise is a direct HDF5 read. Not a
  practical issue but not free either.
\item \textbf{Temperature-interpolation quality.} Both
  representations interpolate between library temperatures
  at runtime. Neither performs on-the-fly Doppler
  broadening from first principles --- WMP is the only
  method of the three in this paper that does. SVD with
  Ducru kernels gives $\sim\!0.1\%$ error on broadened
  cross sections over $300$--$3000$~K, which is tight but
  not analytical.
\item \textbf{Rank dispersion across reactions.} Most
  reactions are rank-one to machine precision ($43$ of
  $47$ non-redundant $^{235}$U channels,
  Sec.~\ref{sec:spectrum}); a minority, chiefly those with
  dense-peak resonance structure in the RRR, need $k \geq
  5$. Using a uniform rank across all reactions wastes
  bytes on the easy channels; per-reaction rank would
  recover those but complicates the dispatch.
\item \textbf{No safe extrapolation outside library
  temperatures.} The Ducru basis is a projection onto the
  training-temperature set. Extrapolating to $T$ outside
  $[T_\text{min}^\text{lib}, T_\text{max}^\text{lib}]$ is
  unsafe in the same way pointwise interpolation is
  unsafe outside its range. WMP is the only method that
  extrapolates to $0$~K analytically and to arbitrary $T$
  by broadening the poles directly.
\end{itemize}

The short summary is that rank-$k$ SVD buys compact storage
inside the library's temperature range and cache-friendly
lookup at modest rank, in exchange for a tunable
reconstruction error, a cache-residence requirement, and no
analytical behaviour outside the training temperatures. The
$132\times$ and $1.06\times$ memory numbers in
Sec.~\ref{sec:hybrid} are both compatible with this picture;
they measure different slices of it.

\subsection{Hybrid SVD+WMP dispatch}

Where $\log \sigma$ is intrinsically high rank
(resolved-resonance region, RRR), we combine SVD with
WMP~\cite{josey2015wmp,forget2014wmp}. WMP represents $\sigma$
analytically via pole--residue pairs arranged into per-energy
windows plus a low-order Doppler-broadened polynomial
background:
\begin{equation}
\sigma_r(E, T) \approx \sum_{p \in \text{window}(E)}
\Re\!\left[
r_{p,r}\,
w\!\left(\frac{\sqrt{E} - p}{\sqrt{kT/\text{AWR}}\cdot 2}\right)
\right] + \mathcal{P}_r(E, T),
\end{equation}
where $w(\cdot)$ is the Faddeeva function
$w(z) = e^{-z^2}\,\mathrm{erfc}(-iz)$ and $\mathcal{P}_r(E, T)$
is the broadened curvefit. We implement $w$ via Humlicek's W4
four-region rational approximation
(\texttt{src/wmp.rs::faddeeva}); OpenMC's sign convention
$w_\text{om}(z) = -\overline{w(\bar z)}$ for
$\mathrm{Im}\,z < 0$ is used so pole residues in conjugate
pairs contribute with correct signs.

A \texttt{HybridSvdWmpXsProvider} wraps the SVD provider. For
each nuclide with WMP coverage and each energy $E \in
[E_\text{min}^\text{WMP}, E_\text{max}^\text{WMP}]$, the
elastic, fission, and capture microscopic cross sections are
taken from WMP; inelastic, $(n,2n)$ and $(n,3n)$ always come
from SVD (WMP does not represent those channels). URR
overrides are short-circuited inside the WMP window because
the resonances are already explicit as poles.

\subsection{Total cross section and charged-particle residue}
\label{sec:total_xs}

OpenMC's HDF5 layout exposes MT$=3$ (total nonelastic) for
most nuclides. When present we assemble $\sigma_t =
\sigma_\text{MT2} + \sigma_\text{MT3}$ exactly; otherwise we
sum non-redundant leaf reactions, excluding sum-MTs
($\{1, 3, 4, 18, 27, 101\}$) and KERMA/heating channels. In
SVD or hybrid mode we set $\sigma_\text{capture} = \max(0,
\sigma_t - \sum_{r \neq \text{cap}} \sigma_r)$ so
charged-particle emission (MT$=103$, $107$, \ldots) and
first-chance fission (MT$=19$, $20$, $21$) are absorbed into
the capture macroscopic rather than dropped. An XS audit
(Sec.~\ref{sec:offset}) confirms this convention matches
OpenMC to $<\!0.25\%$ on total XS across all nine nuclides.

%% ============================================================
